\scanpage{10}
partie ``arithmétique'' -- celle-ci est un sous-groupe ouvert de $\Gamma_{\overline{\mathbb Q}/\mathbb Q}$, qui, on l'a vu, correspond au sous-corps $K_0$ de $\overline{\mathbb Q}/\mathbb Q$, extension finie $/\mathbb Q$ de $\overline{\mathbb Q}/\mathbb Q$, clôture alg de $\mathbb Q$ dans $K$, de sorte qu'on a une suite exacte

\noindent (8)\quad\diagram{\begin{tikzcd}[ampersand replacement=\&,column sep=small,row sep=small]
1\arrow[r,dash]\&\pi_{\bar K/K}\arrow[r]\&E_{\bar K/K}\arrow[r]\&\Gamma_{\bar K/K}\arrow[r]\arrow[d,hook]\&1\\
{}\&{}\&{}\&\Gamma_{\overline{\mathbb Q}/\mathbb Q}\&{}
\end{tikzcd}}

\margin{\uline{\unsure{NB}} $\pi_{\bar K/K}=(1)$ ssi $K$ alg. sur $\mathbb Q$, i.e. fini sur $\mathbb Q$}
On va donner une interprétation de ce noyau, et de la suite exacte (8), en écrivant
\[
(9)\qquad K=\varinjlim_i A_i
\]
où les $A_i$ sont les sous-$\mathbb Q$-algèbres de type fini de $K$, correspondant au syst. projectif des ``modèles affines'' $U_i=\operatorname{Spec}(A_i)$ de $K/\mathbb Q$. Parmi les $A_i$, il y a d'ailleurs un syst. cofinal formé des $A_i$ réguliers, \inserted{i.e. les $U_i$ lisses$/\mathbb Q$,} avec comme morphismes de transition des morphismes de localisation $A_i\longrightarrow A_j=(A_i)_{f_{ij}}$ ($f_{ij}\in A_i\setminus\{0\}$) -- donc les $U_j\unsure{-|\longrightarrow}U_i$ des \unsure{inclusions|immersions} ouvertes, $U_j=(U_i)_{f_{ij}}$. On peut
\margin{9}
